\honest{Joy in the Merely Real}{Eliezer Yudkowsky}{2008}

I quote Keats's ``dull catalogue of common things'' once more, and Feynman: ``Nothing is `mere'.''

What goes into that catalogue? Things that are normal, unmagical, known or knowable, that play by rules at all, which makes them boring: ``in a word, real.'' That is setting yourself up for a fall. Sooner or later everything will disappoint you: either it will not exist, or, worse, it will turn out to be real. ``If we cannot take joy in things that are merely real, our lives will always be empty.''

The rainbow's sin was to have a scientific explanation. Keats writes ``We know her woof,'' an interesting use of ``we.'' I suspect that Keats did not know the explanation himself, that being told someone else knew was too much for him, and that even the idea of an explanation in principle would have been. \nb{The only evidence these posts give about what Keats knew is a dinner-party toast; see ``Fake Reductionism.''} If Keats did not think like that, I know plenty of people who do.

I have said before that nothing is mysterious in itself. Ignorance is a fact about the mind, and to worship a phenomenon for its mystery is to worship your own ignorance. A blank map does not mean a blank territory, only somewhere we have not visited yet. So everything that exists is liable to end up in the dull catalogue.

You can decide that real, explicable things are still worth caring about, or you can suffer an ennui that has no cure. Self-deception is not open to you.

So scientists, who become fascinated by pocket lint, bird droppings and rainbows, are not so strange. ``You might say'' that some of them are the people able, in principle, to enjoy life in the real universe.
